\section{Measuring Autonomous Systems: Toward Reporting Standards}\label{sec:measurement}

Every claim reviewed so far rests on an instrument, and the instruments are often wrong. This section argues that the field's evidence problem is as much one of measurement as of capability, and proposes a minimum reporting standard.

\subsection{The Benchmark Validity Problem}

Audits of widely used benchmarks find errors large enough to reverse conclusions, in both directions. A checklist study found that flaws in task setup or reward design can mis-state agent performance ``by up to 100\% in relative terms''; on $\tau$-bench airline, an agent that returns empty responses scores 38\% and outperforms a GPT-4o-based agent \citep{A18}. Applying the checklist to CVE-Bench reduced overestimation by 33\% in absolute terms, a result the authors obtained by applying their own checklist. Industry audits report the same. OpenAI, which created SWE-bench Verified, audited 138 problems that o3 did not consistently solve and found that 59.4\% ``contained material issues in test design and/or problem description'' \citep{C15}; it recommended SWE-Bench Pro instead, then estimated that about 30\% of Pro tasks are broken (27.4\% flagged by a human-supervised agent review, 34.1\% by human annotators) and retracted that recommendation \citep{V31}. Authors from Scale AI and Reflection AI repaired or retired 480 of the 731 public SWE-Bench Pro tasks (66\%); on the cleaned split, frontier agents resolve up to 99.4\%, and the task explains 32\% of the variance in outcomes against 2.7\% for the model and 0.3\% for the scaffold \citep{C12}. Errors also hide capability: in a data-engineering benchmark, 82.7\% of 81 failed tasks contained at least one benchmark-attributable error, and correcting the answer key raised the same agent and model from 22.66\% to 32.51\% \citep{D13}.

Evaluation code fails as well as task design. An audit of two public leaderboards found that a grader had scored only about half of one agent's episodes and stored the rest as failures, and that a cost column was ``Exactly 2$\times$ in 70 of 134 runs'' because nested calls were counted twice \citep{F12}. Section~\ref{sec:reliability} showed that run-to-run variance is large enough to reorder agents \citep{F19,F24}, so an unrepeated score is itself a measurement error. Practitioner tooling encodes some remedies. One evaluation harness requires each benchmark adapter to score 100\% with its reference (``oracle'') solution before comparison, ships a no-op agent, and counts two implementations as matching only if their run-score ranges overlap, with at least two runs per side (\repofile{harbor-framework/harbor}{adapters/ADAPTER_CONTRIBUTING.md}). Yet its own parity summary still lists merged adapters compared on a single run, including one where Claude Code with Claude Haiku 4.5 scored 43.7 on the original benchmark and 62.6 on the adapted one (\repofile{harbor-framework/harbor}{adapters/parity_summary.csv}). The rule and the record diverge even where the rule is written down.

\subsection{A Minimum Reporting Standard}

Table~\ref{tab:reporting} proposes eight items that a study of autonomous agents should report. Each answers a failure documented in this review. Repeated runs answer variance; a budget-matched single agent answers the cost-matching problem of Section~\ref{sec:architecture}; cost with its spread answers the cost evidence above; a grader audit and task proofs answer broken instruments; held-out tasks answer overfitting, since in one survey of 17 agent benchmarks most lacked an appropriate holdout and seven had none and no stated plan for one \citep{A10}; interventions per task answer the gap between ``autonomous'' and steered work; and labeling the grader answers self-grading, since in one controlled study LLM verifiers approved 49.4\% of submissions that a deterministic grader rejected \citep{T18}. None of the items needs new methods. Most are counts that evaluation pipelines already hold, as the graded-fraction proposal makes plain \citep{F12}.

\begin{table}[t]
\centering
\small
\caption{Proposed minimum reporting standard for evaluations of autonomous agents.}
\label{tab:reporting}
\begin{tabularx}{\linewidth}{@{}p{0.2\linewidth}XX@{}}
\toprule
Item & What to report & Why \\
\midrule
$k$ runs per arm, with spread & Number of independent runs per arm ($k \geq 3$); mean with standard deviation or range; pass@k and pass\textasciicircum k & Single-run scores move 2.2--6.0 points with the run chosen \citep{F19}; on one task, identical runs varied more than agent--model pairings differed \citep{F24} \\
Budget-matched single-agent baseline & One strong agent given the same token or dollar budget as the system under test & In a controlled ablation, a full team was not significantly better than a single agent (mean $+0.5$ points, $p = 0.20$) \citep{T18}; for similar accuracy, cost can differ by almost two orders of magnitude \citep{A10} \\
\$/task with spread & Dollars per task and per passed task (cost-of-pass), from API usage, with spread across runs & Runs on one task differ up to 30-fold in tokens \citep{F6}; cost-of-pass \citep{F8}; double-counted cost columns \citep{F12} \\
Grader audit & How the grader was checked; share of tasks repaired or retired; graded fraction & Broken tests and tasks in major benchmarks \citep{C15,V31,C12}; answer-key errors \citep{D13}; ungraded episodes scored as failures \citep{F12} \\
Task proofs & A reference (oracle) solution scores full marks; an empty run scores zero & An empty-response agent scores 38\% on $\tau$-bench airline \citep{A18}; one harness requires the oracle check before any comparison (see text) \\
Held-out tasks & Tasks unseen during development, scored once & Most surveyed agent benchmarks lack an appropriate holdout \citep{A10} \\
Interventions per task & Human interventions per task and hours run unattended & One lab estimates success without human intervention at 35\% on its 4--8 hour tasks \citep{F32} \\
Independent vs self-graded & Who graded; whether the builder could see or edit the grader & LLM verifiers approved 49.4\% of grader-rejected submissions \citep{T18} \\
\bottomrule
\end{tabularx}
\end{table}

\subsection{How the Field Reports Today}

We did not code the corpus against Table~\ref{tab:reporting}; a systematic count belongs on the research agenda (Section~\ref{sec:agenda}). The sources above suggest the gap is wide. One 2024 analysis found that cost was not a top-line metric in any of the code-generation agent papers it re-evaluated \citep{A10}; the benchmark OpenAI recommended as a replacement was itself found to have broken tasks and the recommendation was retracted \citep{C15,V31}; and even a harness with a written parity rule still lists merged single-run comparisons. For anyone building or buying such systems, the implication is to treat a score that lacks these items as a claim rather than a measurement.
